\documentclass[11pt]{article}
\usepackage[margin=1in]{geometry}
\usepackage{listings}
\usepackage{xcolor}
\usepackage{graphicx}
\usepackage{float}
\usepackage[hidelinks]{hyperref}

\lstset{
    language=C++,
    basicstyle=\ttfamily\small,
    keywordstyle=\color{blue},
    commentstyle=\color{gray},
    stringstyle=\color{red},
    breaklines=true,
    showstringspaces=false,
    tabsize=4,
    numbers=left,
    numberstyle=\tiny,
    frame=tb
}


\title{Laborator 6 - SO}
\author{Rus Lucian-Florin}
\date{May 5, 2026}

\begin{document}

\maketitle

\newpage
\section*{Comenzi laborator}

\begin{figure}[H]
    \centering
    \includegraphics[width=0.8\textwidth]{1.png}
    \caption{Set comenzi laborator 1}
    \label{fig:1}
\end{figure}

In imaginea de mai sus se poate vedea cum redirectionarea comenzii cauzeaza problema de \textbf{buffering}, o problema comuna, cauzata de buffer-ul intern al implementarii \textbf{printf}. Aparitia problemei poate fi prevenita prin utilizarea functiei \textbf{fflush}. Referinta \textbf{\href{https://www.youtube.com/watch?v=tcYo6hipaSA&list=PLfqABt5AS4FkW5mOn2Tn9ZZLLDwA3kZUY&index=2}{video}}.

\newpage
\begin{figure}[H]
    \centering
    \includegraphics[width=0.8\textwidth]{2.png}
    \caption{Set comenzi laborator 2}
    \label{fig:2}
\end{figure}


\begin{figure}[H]
    \centering
    \includegraphics[width=0.8\textwidth]{3.png}
    \caption{Set comenzi laborator 3}
    \label{fig:3}
\end{figure}

\newpage
\section*{Teme}
\subsection*{Problema 2}
\begin{lstlisting}
#include <stdio.h>
#include <string.h>
#include <sys/wait.h>
#include <unistd.h>

int main(void) {
    printf("starting custom shell...\n");

    while (true) {
        printf("\033[1;32muser\033[0m$ ");
        char aux[256] = {0};
        scanf("%s", aux);

        // exit command
        if (strcmp(aux, "exit") == 0) {
            break;
        }

        // fork to execute command
        int pid = fork();

        if (pid == 0) {  // Child process
            execlp(aux, aux, NULL);
            perror("execlp");
            return 1;
        } else {  // Parent process
            wait(NULL);
        }
    }

    printf("stopping custom shell...\n");
    return 0;
}
\end{lstlisting}

\begin{figure}[H]
    \centering
    \includegraphics[width=0.8\textwidth]{4.png}
    \caption{Output \textbf{tellwait2} dupa modificare}
    \label{fig:4}
\end{figure}

\newpage
\subsection*{Problema 3}
\begin{lstlisting}
#include "ourhdr.h"
#include "sys/wait.h"
#include <sys/types.h>

static void charatatime(char *);

int main(void) {
    pid_t pid;

    TELL_WAIT();

    if ((pid = fork()) < 0)
        err_sys("fork error");
    else if (pid == 0) {
        // WAIT_PARENT();		/* parent goes first */
        charatatime("output from child\n");
    }
    else {
        wait(NULL); // folosim `wait` pentru ca e mai usor de folosit
        charatatime("output from parent\n");
        TELL_CHILD(pid);
    }
    exit(0);
}

static void charatatime(char *str) {
    char *ptr;
    int   c;

    setbuf(stdout, NULL); /* set unbuffered */
    for (ptr = str; c = *ptr++;) putc(c, stdout);
}
\end{lstlisting}

\begin{figure}[H]
    \centering
    \includegraphics[width=0.8\textwidth]{5.png}
    \caption{Demonstratie \textbf{nshell}}
    \label{fig:5}
\end{figure}

\textbf{Nota:} comanda \textbf{nshell} suporta doar comenzi formate dintr-un singur cuvant, fara argumente.

\end{document}
